\documentclass[10pt]{article}
\usepackage[margin=0.72in]{geometry}
\usepackage{amsmath,amssymb,booktabs,microtype}
\usepackage[dvipsnames]{xcolor}
\usepackage[hidelinks]{hyperref}
\usepackage{enumitem}

\definecolor{accent}{HTML}{245A73}
\definecolor{soft}{HTML}{F2F6F8}
\setlength{\parindent}{0pt}
\setlength{\parskip}{5pt}
\setlist[itemize]{leftmargin=1.3em,itemsep=2pt,topsep=2pt}
\newcommand{\nrmse}{\mathrm{NRMSE}}
\newcommand{\phina}{\phi_{\mathrm{Na}}}
\newcommand{\phik}{\phi_{\mathrm{K}}}

\title{\vspace{-2.0em}\textcolor{accent}{\bfseries Symbolic validity of numerical NeuronBench recoveries}}
\author{Post-hoc analysis of evolution runs 313195 and 313196}
\date{25 August 2026}

\begin{document}
\maketitle
\vspace{-1.5em}

\colorbox{soft}{\parbox{0.965\linewidth}{
\textbf{Bottom line.} Among 45 held-out fits, 29 attained $\nrmse\le10^{-6}$.
The project's SRBench-style symbolic checker labels 26 of these 29 as ground-truth
matches.  None is literally identical when its fitted decimal constants are treated
as exact numbers, but every recovered frontier contains a polynomial with the correct
ground-truth monomial support.  Four best equations contain additional, numerically
negligible monomials.
}}

\section*{Question and method}

The fully observable target is the current-balance law
\begin{align*}
f = I_{\rm ext}
  &-120\phina(V-50)-36\phik(V+77)-0.3(V+54.4)\\
  &-g_q\phi_q(V-E_q),
\end{align*}
where the final channel depends on the neuron world (and is absent for
\texttt{na\_fatigue}).  PySR fits $f/s$, where $s$ is the training-target RMS.
For this analysis each saved equation was remapped from $x_i$ to the corresponding
physical feature, multiplied by $s$, expanded with SymPy, and compared with the
manifest ground truth.

Three notions must be kept separate:
\begin{itemize}
  \item \textbf{Numerical recovery:} the existing held-out criterion,
  $\nrmse\le10^{-6}$.
  \item \textbf{Project symbolic match:} the same
  \texttt{evaluation.check\_pysr\_symbolic\_match} logic used for SRBench and
  \texttt{evolve\_pysr.py}.  It rounds floating constants to three decimals and
  accepts a zero or constant difference, or a constant ratio.  This is deliberately
  looser than literal algebraic equality.
  \item \textbf{Strict equality/support:} respectively,
  $\operatorname{simplify}(\hat f-f)=0$ with the stored decimals, and equality of
  the monomial sets after polynomial expansion.  An exact target equation must
  itself pass the numerical threshold, so restricting the symbolic sweep to
  recovered frontier rows cannot discard an exact match.
\end{itemize}

\section*{Results}

\begin{center}
\begin{tabular}{lrrrr}
\toprule
Evolution run & Numeric recoveries & Project match & Strict equality & Best support correct\\
\midrule
313196 (trained on one world) & 14 & 14 & 0 & 11\\
313195 (trained on two worlds) & 15 & 12 & 0 & 14\\
\midrule
Total & 29 & 26 & 0 & 25\\
\bottomrule
\end{tabular}
\end{center}

Of the 26 project matches, 23 have zero rounded symbolic difference and three
have a constant rounded difference; none relies only on the constant-ratio rule.
All 29 recovered frontiers contain at least one row with exactly the same monomial
support as the physical ground truth.  Thus the numerical recovery counts mostly
represent recovery of the correct current-balance polynomial, with constants fitted
to finite precision, rather than unrelated interpolating formulas.

\section*{Examples where ``recovered'' does not mean literally identical}

\textbf{1. A tiny extra nonlinear term.}
Run 313196, \texttt{na\_fatigue}, seed 10000 has
$\nrmse=9.06\times10^{-11}$.  Its expanded physical equation differs from ground
truth by
\begin{align*}
\hat f-f={}&-5.54\times10^{-7}\phina^2
 +1.10\times10^{-6}\phina-9.19\times10^{-7}\phik
 -1.07\times10^{-7}\\
&\quad+\text{linear and bilinear corrections of magnitude at most }3.3\times10^{-8}.
\end{align*}
The $\phina^2$ term makes the equation structurally different.  The project checker
still calls it a match because these tiny coefficients disappear under its rounding.

\textbf{2. Correct support, approximate constants, and a project-checker miss.}
Run 313195, \texttt{ca\_rebound}, seed 10000 has
$\nrmse=3.09\times10^{-7}$.  The relevant physical coefficients are:
\begin{center}
\small
\begin{tabular}{lrr@{\hspace{1.2em}}lrr}
\toprule
Term & Ground truth & Found & Term & Ground truth & Found\\
\midrule
$I_{\rm ext}$ & 1 & 0.99999762 & $V\phina$ & -120 & -119.99999734\\
$V\phik$ & -36 & -35.99997305 & $V\phi_T$ & -3.2 & -3.20012909\\
$V$ & -0.3 & -0.29995093 & $\phina$ & 6000 & 5999.99985004\\
$\phik$ & -2772 & -2771.99906434 & $\phi_T$ & 384 & 383.99682787\\
$1$ & -16.32 & -16.31889694 & & &\\
\bottomrule
\end{tabular}
\end{center}
This is the same polynomial support with very close coefficients, yet the project
checker returns no match.  Its three-decimal rounding is applied to the normalized,
factorized PySR expression before expansion, so the decision can depend on how an
otherwise similar polynomial is factored.

\textbf{3. An even smaller checker false negative.}
Run 313195, \texttt{d\_type}, seed 10004 has
$\nrmse=4.30\times10^{-10}$ and the correct support.  For example, the ground-truth
coefficients of $V\phi_D$, $\phi_D$, and the constant are $-9$, $-693$, and
$-16.32$; the fitted values are $-8.999999917$, $-692.999995373$, and
$-16.320003695$.  It also fails the project checker despite being an excellent
coefficient-level recovery.

\section*{Interpretation}

Yes, there are $\nrmse<10^{-6}$ equations that are different from ground truth.
In the strict stored-decimal sense this is true for all 29 recoveries, which is
expected when continuous constants are optimized numerically after target scaling.
The more substantive structural difference occurs in four best equations, all
through tiny extra monomials.  Conversely, three numerically recovered equations
with correct support are missed by the current project checker.

For future NeuronBench summaries, the most informative reporting is therefore a
three-part label: held-out numerical recovery, the existing SRBench-style match,
and canonical expanded-polynomial support (optionally with a stated coefficient
tolerance).  This avoids treating fitted-decimal inequality as a scientific miss
while making small spurious terms visible.

\vfill
\footnotesize
\raggedright
Reproducibility: \path{scripts/analyze_neuron_symbolic_recovery.py} reads
\path{runs/313195/neuron_full_eval/neuron_results.json} and
\path{runs/313196/neuron_full_eval/neuron_results.json}.  Machine-readable
details, including every recovered equation and expanded residual, are stored in
\path{reports/neuron_symbolic_recovery_analysis.json}.

\end{document}
